\chapter{Finite Dimension}\label{ch:b1:findim}

A space spanned by finitely many vectors carries a well-defined
\emph{dimension} --- the common size of all its bases. The proofs
below all flow from one combinatorial engine, the exchange lemma:
a \omterm{def:b1:vspaces:free}{free family} can never outnumber a generating one. With dimension
come the tools used everywhere afterwards: the incomplete \omterm{def:b1:vspaces:free}{basis}
theorem, the \omterm{def:b1:findim:rank}{rank} of a family, Grassmann's formula.

\section{Existence of bases}

\begin{definition}\label{def:b1:findim:def}
$E$ is \emph{finite-dimensional}\index{finite dimension} when it has
a finite generating family. (Otherwise \emph{infinite-dimensional}:
so is $K[X]$, whose finite families only \omterm{def:b1:vspaces:span}{span} \omterm{def:b1:poly:def}{polynomials} of bounded
degree.)
\end{definition}

\begin{theorem}[Exchange lemma]\label{thm:b1:findim:exchange}
Let $(g_1, \dots, g_n)$ generate $E$ and $(f_1, \dots, f_p)$ be \omterm{def:b1:vspaces:free}{free}
in $E$. Then $p \leq n$.
\end{theorem}

\begin{proof}
We prove by induction on $k \leq p$: \emph{after renumbering the
$g_j$, the family $(f_1, \dots, f_k, g_{k+1}, \dots, g_n)$ generates
$E$} --- which forces $k \leq n$ at each stage, and $p \leq n$ at
the end.

$k = 0$: the hypothesis. Step: assume it for $k - 1$; then $f_k$ is
a combination of $(f_1, \dots, f_{k-1},$ $g_k, \dots, g_n)$. In
this combination some $g_j$ ($j \geq k$) has a nonzero coefficient
---
otherwise $f_k$ would be a combination of $f_1, \dots, f_{k-1}$,
contradicting freeness. Renumber so that $j = k$, and solve for
$g_k$: $g_k$ is a combination of $(f_1, \dots, f_k, g_{k+1}, \dots,
g_n)$. Every vector of $E$, expressed through the $(k-1)$-family,
can then be re-expressed through the $k$-family: it generates.
(If $k - 1 = n$, no $g$ remains and $f_k$ would be a combination of
the $f_i$'s alone: impossible; so $k \leq n$.)
\end{proof}

\begin{example}[The exchange, watched once]
\label{ex:b1:findim:exchangerun}
In $\R^2$, take the generating family $\bigl(g_1, g_2\bigr) =
\bigl((1,0), (0,1)\bigr)$ and the \omterm{def:b1:vspaces:free}{free family} $\bigl(f_1,
f_2\bigr) = \bigl((1,2), (3,4)\bigr)$. Step $1$: $f_1 = 1\cdot
g_1 + 2\cdot g_2$; the coefficient of $g_2$ is nonzero, so
exchange $g_2$ for $f_1$: the family $\bigl(f_1, g_1\bigr)$
still generates ($g_2 = \frac12(f_1 - g_1)$). Step $2$: $f_2 =
(3,4) = 2\,f_1 + 1\cdot g_1$; the coefficient of the remaining
$g_1$ is nonzero (it must be: $f_2$ is not a multiple of $f_1$),
so exchange again: $\bigl(f_1, f_2\bigr)$ generates $\R^2$. Had
there been a third \omterm{def:b1:vspaces:free}{free} vector $f_3$, no $g$ would remain to
absorb it --- which is precisely how the lemma forbids $3$ \omterm{def:b1:vspaces:free}{free}
vectors in $\R^2$. The proof above is this bookkeeping done in
general.
\end{example}

\begin{theorem}[Bases in finite dimension]\label{thm:b1:findim:bases}
Let $E \neq \{0\}$ be \omterm{def:b1:findim:def}{finite-dimensional}.
\begin{enumerate}
  \item From every finite generating family one can extract a \omterm{def:b1:vspaces:free}{basis}.
  \item (Incomplete \omterm{def:b1:vspaces:free}{basis} theorem\index{incomplete basis theorem})
        Every \omterm{def:b1:vspaces:free}{free family} extends to a \omterm{def:b1:vspaces:free}{basis}, using vectors of any
        chosen generating family.
  \item All bases of $E$ are finite, with the same number of
        elements: the \emph{dimension} $\dim E$. (Convention:
        $\dim\{0\} = 0$.)
\end{enumerate}
\end{theorem}

\begin{proof}
(1) Discard, one at a time, any vector that is a combination of
the others; the family stays generating, since in any expression
using the discarded vector one may substitute its combination of
the survivors. The process terminates --- each step shrinks a
finite family by one --- and it halts exactly when no remaining
vector is a combination of the others. The final family is still
generating, and it is \omterm{def:b1:vspaces:free}{free}: a nontrivial null combination would
carry some nonzero coefficient, and dividing by it would solve
for the corresponding vector in terms of the others, making it
discardable after all --- contradicting that the process had
halted.

(2) Let $(f_1, \dots, f_p)$ be \omterm{def:b1:vspaces:free}{free}, $(g_1, \dots, g_n)$ generating.
Run through $g_1, \dots, g_n$, appending $g_j$ to the current family
whenever it is not already in its \omterm{def:b1:vspaces:span}{span}
(\cref{prop:b1:vspaces:freecriteria}~(2) keeps the family \omterm{def:b1:vspaces:free}{free}).
The final family is \omterm{def:b1:vspaces:free}{free}, and generating: every $g_j$ lies in its
\omterm{def:b1:vspaces:span}{span} --- either it was appended, or it was already a combination.

(3) Two bases are each \omterm{def:b1:vspaces:free}{free} and each generating: the exchange lemma
gives both inequalities between their cardinalities. (Finiteness:
a \omterm{def:b1:vspaces:free}{basis} is \omterm{def:b1:vspaces:free}{free}, hence by exchange no larger than a finite
generating family.)
\end{proof}

\begin{example}\label{ex:b1:findim:dims}
$\dim K^n = n$ (canonical \omterm{def:b1:vspaces:free}{basis}); $\dim K_n[X] = n + 1$ (monomials);
$\dim_\R \C = 2$; the solution space of $y'' + ay' + by = 0$ has
dimension $2$ (\cref{thm:b1:diffeq:homogeneous2}: the solutions are
parametrized \omterm{def:b1:logic:inj}{bijectively} and linearly by $(\lambda, \mu) \in K^2$).
\end{example}

\begin{example}[The same set, two dimensions]
\label{ex:b1:findim:twodimensions}
The \omterm{def:b1:logic:sets}{set} $\C^2$ of pairs of complex numbers is a $\C$-vector space
of dimension $2$ (canonical \omterm{def:b1:vspaces:free}{basis} $e_1, e_2$) --- and an
$\R$-vector space of dimension $4$, with \omterm{def:b1:vspaces:free}{basis}
\[
(1, 0),\quad (\iu, 0),\quad (0, 1),\quad (0, \iu):
\]
every $(z, w) = (a + \iu b,\ c + \iu d)$ has real \omterm{prop:b1:vspaces:coordinates}{coordinates}
$(a, b, c, d)$, uniquely. Dimension is not a property of the \omterm{def:b1:logic:sets}{set}
of vectors alone: it counts the degrees of freedom \emph{relative
to the allowed scalars}, and halving the scalar supply from $\C$
to $\R$ doubles the count. (The weekend problem exploits the
extreme case of this sensitivity, with scalars shrunk all the
way to $\Q$.)
\end{example}

\begin{example}[Running the completion algorithm]
\label{ex:b1:findim:completion}
Complete the \omterm{def:b1:vspaces:free}{free family} $\bigl((1,1,1)\bigr)$ into a \omterm{def:b1:vspaces:free}{basis} of
$\R^3$ using the canonical vectors. Run the proof of
\cref{thm:b1:findim:bases}~(2) on the generating family $(e_1,
e_2, e_3)$: is $e_1 \in \operatorname{Vect}(1,1,1)$? No (multiples
of $(1,1,1)$ have equal \omterm{prop:b1:vspaces:coordinates}{coordinates}) --- append it. Is $e_2 \in
\operatorname{Vect}\bigl((1,1,1), e_1\bigr)$? A combination
$\alpha(1,1,1) + \beta e_1$ has equal second and third
\omterm{prop:b1:vspaces:coordinates}{coordinates}, and $e_2$ does not --- append it. The family
$\bigl((1,1,1), e_1, e_2\bigr)$ is \omterm{def:b1:vspaces:free}{free} with $3$ vectors: stop,
it is a \omterm{def:b1:vspaces:free}{basis} (\cref{prop:b1:findim:twoofthree} will make this
reflex official). Note the outcome depends on the order in which
the $g_j$ are scanned: completion is an algorithm, not a formula.
\end{example}

\begin{proposition}[The two-out-of-three rule]\label{prop:b1:findim:twoofthree}
Let $\dim E = n$ and $\mathcal{F}$ a family of exactly $n$ vectors
of $E$. Then
\[
\mathcal{F} \text{ free} \iff \mathcal{F} \text{ generating} \iff
\mathcal{F} \text{ basis}.
\]
Moreover any \omterm{def:b1:vspaces:free}{free family} has $\leq n$ vectors, any generating family
$\geq n$.
\end{proposition}

\begin{proof}
The \omterm{def:b1:counting:card}{cardinality} bounds are the exchange lemma against a \omterm{def:b1:vspaces:free}{basis}. If
$\mathcal{F}$ (size $n$) is \omterm{def:b1:vspaces:free}{free} but not generating, some $x$ lies
outside its \omterm{def:b1:vspaces:span}{span}; appending $x$ gives a \omterm{def:b1:vspaces:free}{free family} of $n + 1$
vectors: impossible. If $\mathcal{F}$ is generating but not \omterm{def:b1:vspaces:free}{free},
extracting a \omterm{def:b1:vspaces:free}{basis} (\cref{thm:b1:findim:bases}~(1)) gives a \omterm{def:b1:vspaces:free}{basis} of
$< n$ vectors: impossible.
\end{proof}

\begin{example}[Two-out-of-three, saving half the work]
\label{ex:b1:findim:twoofthreerun}
Is $\bigl((1,1,0), (0,1,1), (1,0,1)\bigr)$ a \omterm{def:b1:vspaces:free}{basis} of $\R^3$?
Count: three vectors, dimension three --- so freeness alone
decides. A null combination gives $a + c = 0$, $a + b = 0$, $b +
c = 0$; adding all three, $2(a + b + c) = 0$, and subtracting
each original equation from $a + b + c = 0$ leaves $b = c = a =
0$: \omterm{def:b1:vspaces:free}{free}, hence a \omterm{def:b1:vspaces:free}{basis}, with the generating half of the
verification supplied by the theorem for free. Compare
\cref{ex:b1:vspaces:testbasis}, where the same double
verification had to be done by hand --- one chapter of theory
converts into exactly that saving, on every \omterm{def:b1:vspaces:free}{basis} check for the
rest of the book.
\end{example}

\begin{method}[Computing a dimension]
\label{met:b1:findim:computedim}
Three standard routes, in decreasing order of frequency.
\begin{enumerate}
  \item \emph{Parametrize, then read off a \omterm{def:b1:vspaces:free}{basis}.} Solve the
        defining constraints, express the general element
        linearly in the surviving parameters, and check that the
        vectors multiplying the parameters are \omterm{def:b1:vspaces:free}{free}: the
        dimension is the number of parameters. (Run below on a
        concrete \omterm{def:b1:vspaces:subspace}{subspace} of $\R^4$.)
  \item \emph{Exhibit a \omterm{def:b1:logic:inj}{bijective} linear parametrization.} When
        the elements are determined by finitely many values ---
        initial conditions of a recurrence
        (\cref{exo:b1:findim:10}), coefficients of a solution
        formula (\cref{ex:b1:findim:dims}) --- the dimension is
        the number of those values.
  \item \emph{Use the formulas.} Grassmann for intersections and
        sums, \omterm{def:b1:findim:rank}{rank} for \omterm{def:b1:vspaces:span}{spans}, and later rank--nullity for
        kernels and images: dimensions are usually
        \emph{computed}, not guessed.
\end{enumerate}
In all three routes, the two-out-of-three rule is the finisher:
once the count matches, freeness or generation alone concludes.
\end{method}

\begin{example}[Route 1, in full]
\label{ex:b1:findim:parametrize}
Dimension of $H = \{(x, y, z, t) \in \R^4 : x + y + z + t = 0
\text{ and } x = t\}$. Solve: $t = x$ and $y + z = -2x$, so $z =
-2x - y$ with $x, y$ \omterm{def:b1:vspaces:free}{free}:
\[
(x,\ y,\ -2x - y,\ x) = x\,(1, 0, -2, 1) + y\,(0, 1, -1, 0) .
\]
The two vectors are \omterm{def:b1:vspaces:free}{free} (look at the first two \omterm{prop:b1:vspaces:coordinates}{coordinates}:
$(x, y) = (0,0)$), so they form a \omterm{def:b1:vspaces:free}{basis} of $H$ and $\dim H = 2$.
The count was predictable --- two independent linear constraints
in $\R^4$ should each eat one dimension --- but the
parametrization proves it \emph{and} hands over a \omterm{def:b1:vspaces:free}{basis}, which
the prediction alone never does; \cref{exo:b1:findim:9} makes
the ``each equation eats at most one dimension'' slogan into a
theorem.
\end{example}

\begin{example}[Route 2, in full]
\label{ex:b1:findim:route2}
Dimension of $W = \{P \in \R_n[X] : P(1) = P(2) = 0\}$ (for $n
\geq 2$). By the factor theorem applied twice
(\cref{thm:b1:poly:factor}; the roots $1$ and $2$ are distinct),
$P \in W$ exactly when $P = (X - 1)(X - 2)\,Q$ with $\deg Q \leq
n - 2$. The correspondence $Q \mapsto (X-1)(X-2)Q$ is linear,
reaches all of $W$, and is \omterm{def:b1:logic:inj}{injective} (a product is zero only if
$Q = 0$): $W$ is parametrized \omterm{def:b1:logic:inj}{bijectively} and linearly by
$\R_{n-2}[X]$, so
\[
\dim W = \dim \R_{n-2}[X] = n - 1 .
\]
A \omterm{def:b1:vspaces:free}{basis} comes with the parametrization: the images of the
monomials, $\bigl((X-1)(X-2),\ (X-1)(X-2)X,\ \dots,\
(X-1)(X-2)X^{n-2}\bigr)$. Each new evaluation constraint at a
fresh point costs exactly one dimension --- the counting
backbone of \omterm{thm:b1:poly:lagrange}{Lagrange interpolation}
(\cref{thm:b1:poly:lagrange}).
\end{example}

\begin{example}[An integral constraint costs one dimension too]
\label{ex:b1:findim:integralconstraint}
Dimension of $H = \{P \in \R_2[X] : \int_0^1 P = 0\}$. Writing
$P = a + bX + cX^2$, the constraint reads $a + \frac b2 + \frac
c3 = 0$; solve for $a$ and parametrize:
\[
P = b\Bigl(X - \frac12\Bigr) + c\Bigl(X^2 - \frac13\Bigr),
\]
so $H = \operatorname{Vect}\bigl(X - \frac12,\ X^2 -
\frac13\bigr)$, dimension $2$ (the two \omterm{def:b1:poly:def}{polynomials} have distinct
degrees: \omterm{def:b1:vspaces:free}{free}). One linear condition --- whether an evaluation,
an \omterm{thm:b1:integration:def}{integral}, or any other linear recipe --- removes at most one
dimension, and exactly one as soon as the condition is not
identically zero. \cref{ch:b1:linmaps} will name such recipes
\emph{linear forms} and their solution \omterm{def:b1:logic:sets}{sets} \emph{hyperplanes};
the \omterm{def:b1:vspaces:free}{basis} vectors found here reappear in
\cref{ch:b1:euclid} as the start of the Legendre family.
\end{example}

\section{Subspaces, rank, Grassmann}

\begin{theorem}[Subspaces]\label{thm:b1:findim:subspaces}
Let $E$ be \omterm{def:b1:findim:def}{finite-dimensional} and $F$ a \omterm{def:b1:vspaces:subspace}{subspace}. Then $F$ is
\omterm{def:b1:findim:def}{finite-dimensional}, $\dim F \leq \dim E$, with equality if and only
if $F = E$. Moreover every \omterm{def:b1:vspaces:subspace}{subspace} has a \omterm{def:b1:vspaces:sum}{supplementary} \omterm{def:b1:vspaces:subspace}{subspace}.
\end{theorem}

\begin{proof}
\omterm{def:b1:vspaces:free}{Free} families of $F$ have at most $\dim E$ vectors (exchange lemma
in $E$: a \omterm{def:b1:vspaces:free}{free family} of $F$ is in particular \omterm{def:b1:vspaces:free}{free} in $E$, and $E$
has a finite generating family). Among the \omterm{def:b1:vspaces:free}{free} families of $F$
pick one of \emph{maximal} size $p$ --- possible since the sizes
are integers bounded by $\dim E$. It generates $F$: otherwise some
$x \in F$ would lie outside its \omterm{def:b1:vspaces:span}{span}, and appending $x$ would give
a \omterm{def:b1:vspaces:free}{free family} of $F$ of size $p + 1$
(\cref{prop:b1:vspaces:freecriteria}~(2)), contradicting
maximality. Being \omterm{def:b1:vspaces:free}{free} and generating, it is a \omterm{def:b1:vspaces:free}{basis} of $F$, and
$\dim F = p \leq \dim E$. If $p = \dim E = n$: a \omterm{def:b1:vspaces:free}{free family} of $n$ vectors of
$E$ is a \omterm{def:b1:vspaces:free}{basis} of $E$ (\cref{prop:b1:findim:twoofthree}), so $F
\supseteq \operatorname{span} = E$. \omterm{def:b1:vspaces:sum}{Supplementary}: complete a \omterm{def:b1:vspaces:free}{basis}
$(f_1, \dots, f_p)$ of $F$ into a \omterm{def:b1:vspaces:free}{basis} $(f_1, \dots, f_p, g_{p+1},
\dots, g_n)$ of $E$ (incomplete \omterm{def:b1:vspaces:free}{basis} theorem); then $G =
\operatorname{Vect}(g_{p+1}, \dots, g_n)$ satisfies $E = F \oplus G$
(existence and uniqueness of decompositions = \omterm{prop:b1:vspaces:coordinates}{coordinates} in the
big \omterm{def:b1:vspaces:free}{basis}).
\end{proof}

\begin{definition}[Rank of a family]\label{def:b1:findim:rank}
The \emph{rank}\index{rank!of a family} of a finite family of
vectors is the dimension of its \omterm{def:b1:vspaces:span}{span}:
$\operatorname{rk}(x_1, \dots, x_p) = \dim
\operatorname{Vect}(x_1, \dots, x_p) \leq \min(p, \dim E)$,
with equality to $p$ iff the family is \omterm{def:b1:vspaces:free}{free}.
\end{definition}

\begin{example}[Computing a rank by elimination]
\label{ex:b1:findim:rankcomputation}
\omterm{def:b1:findim:rank}{Rank} of $\bigl((1,2,3), (2,3,4), (3,4,5), (1,1,1)\bigr)$ in
$\R^3$. The \omterm{def:b1:vspaces:span}{span} is unchanged when one subtracts from a vector a
combination of the others (both families \omterm{def:b1:vspaces:span}{span} the same
combinations): replace $(2,3,4)$ by $(2,3,4) - (1,2,3) = (1,1,1)$
and $(3,4,5)$ by $(3,4,5) - (1,2,3) = (2,2,2)$. The \omterm{def:b1:vspaces:span}{span} is now
$\operatorname{Vect}\bigl((1,2,3), (1,1,1), (2,2,2),
(1,1,1)\bigr) = \operatorname{Vect}\bigl((1,2,3), (1,1,1)\bigr)$,
and these two vectors are not proportional: the \omterm{def:b1:findim:rank}{rank} is $2$. This
subtract-and-discard procedure is systematized as Gaussian
elimination in \cref{ch:b1:det}.
\end{example}

\begin{example}[Summing by concatenation]
\label{ex:b1:findim:sumconcat}
Take, in $\R^3$,
\[
F = \operatorname{Vect}\bigl((1,2,3),\ (1,1,1)\bigr),
\qquad
G = \operatorname{Vect}\bigl((2,3,4)\bigr) .
\]
The sum $F + G$ is spanned by the concatenated family of all
three generators, and
\[
(2, 3, 4) = (1, 2, 3) + (1, 1, 1)
\]
shows the third is redundant: $F + G = F$, of dimension $2$ ---
equivalently $G \subseteq F$, which the relation displays.
Grassmann confirms: $\dim(F \cap G) = 2 + 1 - 2 = 1 = \dim G$.
Sums are computed by concatenating generators and then reducing
the pile by the \omterm{def:b1:findim:rank}{rank} algorithm; no new technique is ever needed.
\end{example}

\begin{theorem}[Grassmann's formula]\label{thm:b1:findim:grassmann}
For \omterm{def:b1:findim:def}{finite-dimensional} \omterm{def:b1:vspaces:subspace}{subspaces} $F, G$ of
$E$:\index{Grassmann's formula}
\[
\dim(F + G) = \dim F + \dim G - \dim(F \cap G) .
\]
In particular $F + G$ is direct iff $\dim(F + G) = \dim F + \dim G$.
\end{theorem}

\begin{proof}
Start from a \omterm{def:b1:vspaces:free}{basis} $(e_1, \dots, e_r)$ of $F \cap G$; complete it
into a \omterm{def:b1:vspaces:free}{basis} $(e_1, \dots, e_r, f_1, \dots, f_s)$ of $F$ and into a
\omterm{def:b1:vspaces:free}{basis} $(e_1, \dots, e_r, g_1, \dots, g_t)$ of $G$
(\cref{thm:b1:findim:bases}~(2)). We claim
\[
\mathcal{B} = (e_1, \dots, e_r, f_1, \dots, f_s, g_1, \dots, g_t)
\]
is a \omterm{def:b1:vspaces:free}{basis} of $F + G$; the formula follows by counting: $(r + s) +
(r + t) - r = r + s + t$.

$\mathcal{B}$ generates $F + G$: any $u + v$ ($u \in F$, $v \in G$)
expands through it. Freeness: suppose $\sum \alpha_i e_i + \sum
\beta_j f_j + \sum \gamma_k g_k = 0$. The vector $w = \sum \gamma_k
g_k = -\sum\alpha_i e_i - \sum\beta_j f_j$ lies in $G \cap F$, so it
expands on $(e_i)$ alone; but $w$ also expands on $(g_k)$, and in
the \omterm{def:b1:vspaces:free}{basis} of $G$ these two expressions must coincide: all
$\gamma_k = 0$ (and the $e$-coordinates match). The relation reduces
to $\sum\alpha_i e_i + \sum\beta_j f_j = 0$, a relation in the \omterm{def:b1:vspaces:free}{basis}
of $F$: all remaining coefficients vanish.
\end{proof}

\begin{example}\label{ex:b1:findim:planes}
Two distinct planes $F, G$ (dimension $2$) of $\R^3$ satisfy $F + G
= \R^3$ (their sum strictly contains a plane), so $\dim(F \cap G) =
2 + 2 - 3 = 1$: they always intersect along a line --- no
``parallel planes'' through the origin.
\end{example}

\begin{example}[Grassmann in action, in $\R^4$]
\label{ex:b1:findim:grassmannrun}
Let $F = \operatorname{Vect}\bigl(e_1,\ e_2,\ (1,1,1,0)\bigr)$
and $G = \operatorname{Vect}(e_3, e_4)$ in $\R^4$. Dimensions:
$\dim F = 3$ (the third generator has a nonzero third coordinate,
outside $\operatorname{Vect}(e_1, e_2)$) and $\dim G = 2$. Sum:
$F + G$ contains $e_1, e_2, e_4$ and $e_3 = (1,1,1,0) - e_1 -
e_2$: it is all of $\R^4$. Grassmann then \emph{computes} the
intersection's size with no elimination at all:
\[
\dim(F \cap G) = 3 + 2 - 4 = 1 .
\]
To identify the line, look inside $G$: a vector $(0, 0, c, d)$
lies in $F$ exactly when it is $a e_1 + b e_2 + \lambda(1,1,1,0)$,
forcing $\lambda = -a = -b$ and $d = 0$: the intersection is
$\operatorname{Vect}(e_3)$ --- consistent, since $e_3$ was
exhibited in $F$ above and lies in $G$ by definition. Typical
division of labor: Grassmann predicts \emph{how much} to look
for, the linear system then finds \emph{what}.
\end{example}

\begin{example}[Intersecting by equations]
\label{ex:b1:findim:intersectequations}
When both \omterm{def:b1:vspaces:subspace}{subspaces} come as solution \omterm{def:b1:logic:sets}{sets}, intersecting is just
stacking the equations. In $\R^3$: $F = \{x + y + z = 0\}$ and
$G = \{x = y\}$ give
\[
F \cap G = \{x = y,\ 2x + z = 0\}
= \{(x,\ x,\ -2x)\} = \operatorname{Vect}(1, 1, -2),
\]
a line. Cross-check by Grassmann: $F + G = \R^3$ (the planes are
distinct, so their sum strictly contains a plane), hence
$\dim(F\cap G) = 2 + 2 - 3 = 1$. The two descriptions of a
\omterm{def:b1:vspaces:subspace}{subspace} --- by equations, by generators --- each make one
operation trivial: equations intersect by stacking, generators
sum by concatenating; converting between them is exactly what
solving a linear system means (\cref{ch:b1:det}).
\end{example}

\begin{remark}[Common pitfalls]
\label{rem:b1:findim:pitfalls}
\emph{Dimensions do not add along sums} unless the sum is direct:
two planes of $\R^3$ have $\dim(F + G) = 3$, not $4$; always
correct by the intersection term (Grassmann).
\emph{An inclusion needs the dimension \emph{and} the inclusion}:
$\dim F = \dim G$ alone never gives $F = G$ (two distinct lines
of $\R^2$); the equality case of
\cref{thm:b1:findim:subspaces} requires $F \subseteq G$ first.
\emph{Counting parameters is not yet a proof}: ``two equations in
$\R^4$, so dimension $2$'' fails when the equations are
dependent ($x + y = 0$ and $2x + 2y = 0$ leave dimension $3$);
only an actual parametrization or a \omterm{def:b1:findim:rank}{rank} computation decides.
\emph{Do not speak of the dimension of a \omterm{def:b1:vspaces:subspace}{non-subspace}}: solution
\omterm{def:b1:logic:sets}{sets} of \emph{inhomogeneous} systems miss $0$; their ``dimension''
is that of the associated homogeneous solution space
(\cref{ch:b1:det} makes this precise).
\emph{Infinite dimension exists}: $K[X]$ contains \omterm{def:b1:vspaces:free}{free} families
of every size (the monomials), so no finite generating family
can exist --- \omterm{def:b1:logic:statement}{statements} like the two-out-of-three rule are
strictly \omterm{def:b1:findim:def}{finite-dimensional} and fail badly on $K[X]$
(\cref{cor:b1:linmaps:samedim} will show the same for
injectivity/surjectivity).
\emph{\omterm{def:b1:findim:rank}{Rank} is about the \omterm{def:b1:vspaces:span}{span}, not the list}: repeating a vector,
reordering, or rescaling by nonzero constants leaves the \omterm{def:b1:findim:rank}{rank}
unchanged, and $\operatorname{rk} = p$ (the number of vectors)
is a \emph{property to prove} --- it is exactly freeness. A
family of $5$ vectors of \omterm{def:b1:findim:rank}{rank} $2$ carries three vectors' worth
of redundancy, which elimination
(\cref{ex:b1:findim:rankcomputation}) locates explicitly.
\end{remark}

\begin{remark}[Where dimension goes to work]
\label{rem:b1:findim:whereused}
Dimension is the book's favorite counting argument from here on.
\cref{ch:b1:linmaps} proves the rank--nullity theorem, the
functional version of Grassmann's formula; \cref{ch:b1:matrices}
computes \omterm{def:b1:findim:rank}{ranks} by row reduction; \cref{ch:b1:det} turns ``$n$
vectors of $K^n$ form a \omterm{def:b1:vspaces:free}{basis}'' into one number being nonzero. The
weekend problem below shows dimension doing \emph{arithmetic}:
counting dimensions over the \omterm{def:b1:structures:field}{field} $\Q$ proves irrationality
\omterm{def:b1:logic:statement}{statements} that look untouchable by hand. In the Year~3 volume the
same dimension counts, refined by \omterm{def:b1:structures:group}{group} theory, decide which
classical construction problems are solvable --- that story is
Galois theory.
\end{remark}

\begin{remark}[Perspectives inside Book 3]
\label{rem:b1:findim:perspectives}
Dimension is the conserved quantity of the rest of this volume,
and it is worth naming the conservation laws in advance.
\cref{ch:b1:linmaps} proves $\dim E = \dim\ker u +
\operatorname{rk} u$: what a linear \omterm{def:b1:logic:map}{map} crushes plus what it
keeps always totals the source. \cref{ch:b1:det} refines this
into the structure of solution \omterm{def:b1:logic:sets}{sets}: $p$ unknowns minus
$\operatorname{rk} A$ pivots leaves the dimension of the
solution space, which Gaussian elimination exhibits as \omterm{def:b1:vspaces:free}{free}
parameters. \cref{ch:b1:euclid} splits $\dim E = \dim F + \dim
F^\perp$ orthogonally, and the weekend problem of
\cref{ch:b1:multivar} spends exactly this budget: $n$ data
points, $2$ parameters fitted, $n - 2$ dimensions of residual.
Whenever a count refuses to balance in a later chapter, the
error is a forgotten kernel or a non-direct sum --- come back to
Grassmann's formula first.
\end{remark}

\section{Exercises}

\begin{exercise}[$\star$]\label{exo:b1:findim:1}
Give a \omterm{def:b1:vspaces:free}{basis} and the dimension of:
\begin{enumerate}
  \item $F = \{(x,y,z) \in \R^3 : x + y + z = 0\}$;
  \item $G = \{(x,y,z,t) \in \R^4 : x = y,\ z = 2t\}$;
  \item $H = \{P \in \R_3[X] : P(1) = P'(1) = 0\}$.
\end{enumerate}
\end{exercise}

\begin{exercise}[$\star$]\label{exo:b1:findim:2}
Compute the \omterm{def:b1:findim:rank}{rank} of the family
$\bigl((1,1,1), (1,2,3), (3,5,7), (0,1,2)\bigr)$ in $\R^3$, and
extract a \omterm{def:b1:vspaces:free}{basis} of its \omterm{def:b1:vspaces:span}{span}.
\end{exercise}

\begin{exercise}[$\star$]\label{exo:b1:findim:3}
Complete the \omterm{def:b1:vspaces:free}{free family} $\bigl((1,1,0,0), (0,0,1,1)\bigr)$ into a
\omterm{def:b1:vspaces:free}{basis} of $\R^4$ using canonical vectors, and justify.
\end{exercise}

\begin{exercise}[$\star$]\label{exo:b1:findim:4}
Prove that $\bigl(1, X, X(X-1), X(X-1)(X-2)\bigr)$ is a \omterm{def:b1:vspaces:free}{basis} of
$\R_3[X]$, and find the \omterm{prop:b1:vspaces:coordinates}{coordinates} of $X^3$ in it.
\end{exercise}

\begin{exercise}[$\star\star$]\label{exo:b1:findim:5}
Let $F$ and $G$ be \omterm{def:b1:vspaces:subspace}{subspaces} of dimensions $4$ and $5$ of a space
$E$ with $\dim E = 7$. What are the possible values of $\dim(F \cap
G)$? Give an instance realizing each value with $E = \R^7$.
\end{exercise}

\begin{exercise}[$\star\star$]\label{exo:b1:findim:6}
Let $H = \{P \in \R_n[X] : P(1) = 0\}$. Prove that $H$ is a
hyperplane of $\R_n[X]$ (a \omterm{def:b1:vspaces:subspace}{subspace} of dimension $n$), exhibit a
\omterm{def:b1:vspaces:free}{basis} of $H$ \emph{(think of the factor theorem: $P = (X-1)Q$)}, and
give a \omterm{def:b1:vspaces:sum}{supplementary} line.
\end{exercise}

\begin{exercise}[$\star\star$]\label{exo:b1:findim:7}
Let $u_1, \dots, u_p$ be vectors of \omterm{def:b1:findim:rank}{rank} $r$. Prove that removing
one vector yields a family of \omterm{def:b1:findim:rank}{rank} $r$ or $r - 1$, and that
appending one vector yields \omterm{def:b1:findim:rank}{rank} $r$ or $r + 1$. Deduce that \omterm{def:b1:findim:rank}{rank}
changes by at most $1$ under any single insertion or deletion.
\end{exercise}

\begin{exercise}[$\star\star\star$]\label{exo:b1:findim:8}
Let $F_1 \subseteq F_2 \subseteq \dots \subseteq F_k$ be \omterm{def:b1:vspaces:subspace}{subspaces}
of $E$ ($\dim E = n$) with $F_i \neq F_{i+1}$ for all $i$. Prove $k
\leq n + 1$. Deduce that a strictly increasing chain of \omterm{def:b1:vspaces:subspace}{subspaces}
of $\R^n$ has length at most $n + 1$, and exhibit one of maximal
length.
\end{exercise}

\begin{exercise}[$\star\star\star$]\label{exo:b1:findim:9}
Let $E$ be of dimension $n$ and $F$, $G$ two hyperplanes (dimension
$n - 1$), $F \neq G$. Compute $\dim(F \cap G)$. Generalize: the
intersection of $k$ hyperplanes has dimension $\geq n - k$.
\end{exercise}

\begin{exercise}[$\star\star$]\label{exo:b1:findim:10}
Let $E$ be the \omterm{def:b1:logic:sets}{set} of real sequences satisfying $u_{n+2} =
3u_{n+1} - 2u_n$ for all $n$.
\begin{enumerate}
  \item Show that $E$ is a \omterm{def:b1:vspaces:subspace}{subspace} of the space of sequences, and
        that a sequence of $E$ is entirely determined, linearly,
        by the pair $(u_0, u_1)$; deduce $\dim E = 2$.
  \item Check that the constant sequence $(1)$ and the geometric
        sequence $(2^n)$ lie in $E$ and form a \omterm{def:b1:vspaces:free}{basis} of $E$.
  \item Find the sequence of $E$ with $u_0 = 0$, $u_1 = 1$.
\end{enumerate}
\end{exercise}

\begin{exercise}[$\star\star$]\label{exo:b1:findim:11}
Let $E$ be of dimension $n$.
\begin{enumerate}
  \item If $F, G$ are \omterm{def:b1:vspaces:subspace}{subspaces} with $\dim F + \dim G > n$, prove
        $F \cap G \neq \{0\}$. Illustrate: two \omterm{def:b1:vspaces:subspace}{subspaces} of
        dimensions $51$ and $50$ of $\R^{100}$ always share a
        nonzero vector.
  \item If $H$ is a hyperplane and $F$ a \omterm{def:b1:vspaces:subspace}{subspace} with $F \cap H =
        \{0\}$, prove $\dim F \leq 1$.
\end{enumerate}
\end{exercise}

\begin{exercise}[$\star\star\star$]\label{exo:b1:findim:12}
(Common \omterm{def:b1:vspaces:sum}{supplementary}) Let $F, G$ be \omterm{def:b1:vspaces:subspace}{subspaces} of $E$ (\omterm{def:b1:findim:def}{finite
dimension}) with $\dim F = \dim G$. Prove that $F$ and $G$ admit a
\emph{common} \omterm{def:b1:vspaces:sum}{supplementary}: there is a \omterm{def:b1:vspaces:subspace}{subspace} $S$ with $E = F
\oplus S = G \oplus S$. \emph{(Induct downward on $\dim F$: if $F
\neq G \neq E$, pick $x \notin F \cup G$ ---
\cref{exo:b1:vspaces:12} allows it --- and consider $F \oplus
\operatorname{Vect}(x)$ and $G \oplus \operatorname{Vect}(x)$.)}
\end{exercise}

\section{Problem: Dedekind's tower law}

\begin{problem}\label{pb:b1:findim:1}
Nothing in Chapters 18--19 used anything about the scalars beyond
the \omterm{def:b1:structures:field}{field} axioms (\cref{def:b1:structures:field}): one may
therefore take $K = \Q$ and measure \omterm{def:b1:logic:sets}{sets} of \emph{real numbers}
with the yardstick of $\Q$-dimension. This problem computes the
dimension of $\Q(\sqrt2, \sqrt3)$, proves \emph{Dedekind's \omterm{pb:b1:findim:1}{tower
law}} $\dim_\Q M = \dim_\Q K \cdot \dim_K M$, and harvests
irrationality theorems by pure dimension counting --- no
$\varepsilon$, no decimals, just bases.\index{tower law}
\index{field!subfield of $\R$}\index{irrational number}

\medskip
\noindent\textbf{Part I --- Rational scalars.}
\begin{enumerate}
  \item Check that $\Q$ is a \omterm{def:b1:structures:field}{field} and that every definition and
        proof of Chapters 18--19 uses only the \omterm{def:b1:structures:field}{field} axioms of the
        scalars; conclude that $\R$ is a $\Q$-vector space and
        that the exchange lemma, the \omterm{def:b1:vspaces:free}{basis} theorems and
        Grassmann's formula hold over $\Q$. Point out the one step
        of the proof of \cref{thm:b1:findim:exchange} where
        division by a nonzero scalar is performed.
  \item Show that $(1, \sqrt2)$ is \omterm{def:b1:vspaces:free}{free} over $\Q$ but linked over
        $\R$. Set $\Q(\sqrt2) = \operatorname{Vect}_\Q(1,
        \sqrt2) = \{a + b\sqrt2 : a, b \in \Q\}$; what is
        $\dim_\Q \Q(\sqrt2)$?
  \item Show that $\Q(\sqrt2)$ is stable under multiplication.
        For $u = a + b\sqrt2$ set $\sigma(u) = a - b\sqrt2$ and
        $N(u) = u\,\sigma(u) = a^2 - 2b^2$. Show $\sigma(uv) =
        \sigma(u)\sigma(v)$, deduce $N(uv) = N(u)N(v)$, and show
        $N(u) \neq 0$ whenever $u \neq 0$.
  \item Deduce that every nonzero $u \in \Q(\sqrt2)$ has its
        inverse in $\Q(\sqrt2)$, namely $u^{-1} =
        \sigma(u)/N(u)$: $\Q(\sqrt2)$ is a subfield of $\R$.
        Compute $\dfrac1{3 + 2\sqrt2}$ and $\dfrac1{1 + \sqrt2}$.
\end{enumerate}

\medskip
\noindent\textbf{Part II --- Adjoining $\sqrt3$.}
\begin{enumerate}[resume]
  \item Prove that $\sqrt6 \notin \Q$ \emph{(compare the exponent
        of $2$ on both sides of $6q^2 = p^2$, as in
        \cref{exo:b1:arith:7})}, then that $\sqrt3 \notin
        \Q(\sqrt2)$ \emph{(square $\sqrt3 = a + b\sqrt2$ and
        discuss the cases $ab \neq 0$, $b = 0$, $a = 0$)}.
  \item Let $M = \operatorname{Vect}_\Q(1, \sqrt2, \sqrt3,
        \sqrt6)$. Show that $M$ is stable under multiplication
        \emph{(a table of the products of \omterm{def:b1:vspaces:free}{basis} vectors
        suffices)}.
  \item Write $K = \Q(\sqrt2)$. Show that $(1, \sqrt3)$ is \omterm{def:b1:vspaces:free}{free}
        \emph{over the \omterm{def:b1:structures:field}{field} $K$}, and deduce that $M = K +
        K\sqrt3$ is a $K$-vector space of dimension $2$ with
        \omterm{def:b1:vspaces:free}{basis} $(1, \sqrt3)$.
  \item Prove that $(1, \sqrt2, \sqrt3, \sqrt6)$ is \omterm{def:b1:vspaces:free}{free} over
        $\Q$, hence $\dim_\Q M = 4$. \emph{(Group a null relation
        as $(a + b\sqrt2) + (c + d\sqrt2)\sqrt3 = 0$ and apply
        questions 7 then 2.)}
\end{enumerate}

\medskip
\noindent\textbf{Part III --- The \omterm{pb:b1:findim:1}{tower law}.} Let $\Q \subseteq K
\subseteq M \subseteq \R$ where $K$ and $M$ are subfields, $(e_1,
\dots, e_m)$ is a \omterm{def:b1:vspaces:free}{basis} of $K$ as a $\Q$-vector space, and $(f_1,
\dots, f_n)$ a \omterm{def:b1:vspaces:free}{basis} of $M$ as a $K$-vector space.
\begin{enumerate}[resume]
  \item Show that the $mn$ products $(e_i f_j)$ generate $M$ over
        $\Q$.
  \item Show that the family $(e_i f_j)$ is \omterm{def:b1:vspaces:free}{free} over $\Q$.
        \emph{(Reorganize a null $\Q$-combination as $\sum_j
        \bigl(\sum_i \lambda_{ij} e_i\bigr) f_j$, whose inner
        coefficients live in $K$.)}
  \item Conclude with \emph{Dedekind's \omterm{pb:b1:findim:1}{tower law}}:
        \[
        \dim_\Q M \;=\; \dim_\Q K \,\cdot\, \dim_K M ,
        \]
        and check it on $\Q \subseteq \Q(\sqrt2) \subseteq M$
        against questions 7 and 8.
  \item (\omterm{def:b1:structures:field}{Fields} for free) Let $A \subseteq \R$ be a
        \omterm{def:b1:findim:def}{finite-dimensional} $\Q$-subspace containing $1$ and
        stable under multiplication, and let $u \in A$, $u \neq
        0$. Show that if $(a_1, \dots, a_d)$ is a \omterm{def:b1:vspaces:free}{basis} of $A$,
        then $(u a_1, \dots, u a_d)$ is again a \omterm{def:b1:vspaces:free}{basis} of $A$;
        deduce that $u$ has an inverse \emph{in $A$}: $A$ is a
        subfield of $\R$. Which earlier questions does this
        recover?
  \item Show that for every $x \in A$ (as in question 12, $\dim_\Q
        A = d$) the family $(1, x, x^2, \dots, x^{d})$ is linked:
        every element of $A$ is a root of a nonzero \omterm{def:b1:poly:def}{polynomial}
        with rational coefficients, of degree at most $d$.
\end{enumerate}

\medskip
\noindent\textbf{Part IV --- One number generates everything.}
Set $s = \sqrt2 + \sqrt3$.
\begin{enumerate}[resume]
  \item Compute the \omterm{prop:b1:vspaces:coordinates}{coordinates} of $s^2$, $s^3$ and $s^4$ in the
        \omterm{def:b1:vspaces:free}{basis} $(1, \sqrt2, \sqrt3, \sqrt6)$ of $M$.
  \item Show that $(1, s, s^2, s^3)$ is \omterm{def:b1:vspaces:free}{free} over $\Q$, and deduce
        $\operatorname{Vect}_\Q(1, s, s^2, s^3) = M$: every
        element of $M$ is a rational \omterm{def:b1:poly:def}{polynomial} in $s$. Express
        $\sqrt2$ and $\sqrt3$ as such \omterm{def:b1:poly:def}{polynomials}.
  \item Verify $s^4 - 10s^2 + 1 = 0$, and show that $X^4 - 10X^2
        + 1$ is the \emph{\omterm{def:b1:poly:def}{monic polynomial} of least degree}
        vanishing at $s$. Determine its four real roots.
  \item Deduce: $1/s = 10s - s^3$; identify this number. Show
        that $a + b\sqrt2 + c\sqrt3 + d\sqrt6$ ($a,b,c,d \in \Q$)
        is rational if and only if $b = c = d = 0$; in
        particular $\sqrt2 + \sqrt3 + \sqrt6$ is irrational
        (strengthening \cref{exo:b1:reals:7}).
\end{enumerate}

\medskip
\noindent\textbf{Part V --- The cube root stays outside.} Set $t
= 2^{1/3}$.
\begin{enumerate}[resume]
  \item Show that $X^3 - 2$ has no rational root \emph{(the
        rational-root criterion of \cref{exo:b1:poly:5}, or a
        valuation count)}; deduce that $(1, t)$ is \omterm{def:b1:vspaces:free}{free} over
        $\Q$.
  \item Show that $(1, t, t^2)$ is \omterm{def:b1:vspaces:free}{free} over $\Q$. \emph{(If some
        nonzero $P \in \Q[X]$ of degree $\leq 2$ kills $t$, take
        one of least degree and divide $X^3 - 2$ by it,
        \cref{thm:b1:poly:division}; conclude that $X^3 - 2$
        would have a rational root.)} Deduce that $A =
        \operatorname{Vect}_\Q(1, t, t^2)$ has dimension $3$, is
        stable under multiplication, and is a subfield of $\R$.
  \item Prove that $t \notin M$: otherwise $M$ would be a \omterm{def:b1:vspaces:def}{vector
        space} over the \omterm{def:b1:structures:field}{field} $A$, and the \omterm{pb:b1:findim:1}{tower law} would force
        $3 \mid 4$. So $2^{1/3}$ is not a rational combination
        of $1, \sqrt2, \sqrt3, \sqrt6$.
  \item Deduce that $t \notin \Q(\sqrt2)$, that $t + \sqrt2$ is
        irrational, and that no rationals $a, b$ satisfy
        $2^{1/3} = a + b\sqrt3$.
\end{enumerate}

\medskip
\noindent\textbf{Part VI --- All the subfields, and synthesis.}
\begin{enumerate}[resume]
  \item Prove the \omterm{def:b1:arith:divides}{divisibility} of degrees: if $\Q \subseteq A
        \subseteq B \subseteq \R$ are subfields with $\dim_\Q B$
        finite, then $\dim_\Q A$ \omterm{def:b1:arith:divides}{divides} $\dim_\Q B$. What are
        the possible dimensions of subfields of $M$?
  \item Determine \emph{all} subfields of $M$ of dimension $2$
        over $\Q$. \emph{(Reduce to finding the $w = a + b\sqrt2
        + c\sqrt3 + d\sqrt6 \in M \setminus \Q$ with $w^2 \in
        \Q$: expand $w^2$ and annihilate the three irrational
        \omterm{prop:b1:vspaces:coordinates}{coordinates}.)}
  \item Express $\dfrac1{1 + \sqrt2 + \sqrt3}$ in the \omterm{def:b1:vspaces:free}{basis} $(1,
        \sqrt2, \sqrt3, \sqrt6)$ \emph{(multiply by well-chosen
        conjugates)}.
  \item Synthesis, in four sentences: why $\Q$-dimension is an
        \emph{arithmetic} invariant of a \omterm{def:b1:logic:sets}{set} of reals; what
        finiteness of the dimension forces about every element
        (question 13); how the two-out-of-three rule produced
        inverses out of thin air (question 12); and what the
        \omterm{pb:b1:findim:1}{tower law} forbids (question 20). Name the theorem proved
        in Part III.
\end{enumerate}
\end{problem}
